% Options for packages loaded elsewhere
\PassOptionsToPackage{unicode}{hyperref}
\PassOptionsToPackage{hyphens}{url}
\documentclass[
  11pt,
  a4paper]{article}
\usepackage{xcolor}
\usepackage{amsmath,amssymb}
\setcounter{secnumdepth}{5}
\usepackage{iftex}
\ifPDFTeX
  \usepackage[T1]{fontenc}
  \usepackage[utf8]{inputenc}
  \usepackage{textcomp} % provide euro and other symbols
\else % if luatex or xetex
  \usepackage{unicode-math} % this also loads fontspec
  \defaultfontfeatures{Scale=MatchLowercase}
  \defaultfontfeatures[\rmfamily]{Ligatures=TeX,Scale=1}
\fi
\usepackage{lmodern}
\ifPDFTeX\else
  % xetex/luatex font selection
  \setmainfont[]{Latin Modern Roman}
\fi
% Use upquote if available, for straight quotes in verbatim environments
\IfFileExists{upquote.sty}{\usepackage{upquote}}{}
\IfFileExists{microtype.sty}{% use microtype if available
  \usepackage[]{microtype}
  \UseMicrotypeSet[protrusion]{basicmath} % disable protrusion for tt fonts
}{}
\makeatletter
\@ifundefined{KOMAClassName}{% if non-KOMA class
  \IfFileExists{parskip.sty}{%
    \usepackage{parskip}
  }{% else
    \setlength{\parindent}{0pt}
    \setlength{\parskip}{6pt plus 2pt minus 1pt}}
}{% if KOMA class
  \KOMAoptions{parskip=half}}
\makeatother
\usepackage{longtable,booktabs,array}
\usepackage{calc} % for calculating minipage widths
% Correct order of tables after \paragraph or \subparagraph
\usepackage{etoolbox}
\makeatletter
\patchcmd\longtable{\par}{\if@noskipsec\mbox{}\fi\par}{}{}
\makeatother
% Allow footnotes in longtable head/foot
\IfFileExists{footnotehyper.sty}{\usepackage{footnotehyper}}{\usepackage{footnote}}
\makesavenoteenv{longtable}
\setlength{\emergencystretch}{3em} % prevent overfull lines
\providecommand{\tightlist}{%
  \setlength{\itemsep}{0pt}\setlength{\parskip}{0pt}}

\usepackage{amsmath,amssymb,amsthm}
\usepackage{booktabs}
\usepackage{graphicx}
\usepackage[margin=1in]{geometry}
\usepackage{hyperref}
\usepackage{xcolor}
\usepackage{unicode-math}
\hypersetup{colorlinks=true,linkcolor=blue!60!black,citecolor=green!50!black,urlcolor=blue!60!black}
\usepackage{fancyhdr}
\pagestyle{fancy}
\fancyhf{}
\fancyhead[L]{\small PACSeries Paper \thepapernumber}
\fancyhead[R]{\small Dawn Field Institute}
\fancyfoot[C]{\thepage}
\renewcommand{\headrulewidth}{0.4pt}

% Paper number counter
\newcounter{papernumber}

\setcounter{papernumber}{5}
\usepackage{bookmark}
\IfFileExists{xurl.sty}{\usepackage{xurl}}{} % add URL line breaks if available
\urlstyle{same}
\hypersetup{
  hidelinks,
  pdfcreator={LaTeX via pandoc}}

\title{Classical Physics from Information Geometry}
\author{Peter Groom}
\date{February 2026}

\begin{document}
\maketitle

\renewcommand*\contentsname{Contents}
{
\setcounter{tocdepth}{3}
\tableofcontents
}
\section{Classical Physics from Information
Geometry}\label{classical-physics-from-information-geometry}

\textbf{Peter Groom, Dawn Field Institute}\\
\textbf{PACSeries Paper 5}\\
\textbf{Date}: February 2026\\
\textbf{Version}: 2.1

\begin{center}\rule{0.5\linewidth}{0.5pt}\end{center}

\subsection{Abstract}\label{abstract}

Maxwell's equations are among the most successful structures in physics.
They unify electricity, magnetism, and optics; they predict
electromagnetic waves at speed \(c\); they are Lorentz-invariant before
general relativity was formulated. This paper examines whether their
structure --- curl operations, inverse-square forces, three spatial
dimensions, and the specific value of the fine structure constant ---
can be understood as consequences of information-theoretic constraints
rather than independent empirical postulations.

We show that the SEC wave equation
\(\partial^2 S/\partial t^2 = (c_\alpha c_\gamma + c_\beta c_\delta)\nabla^2 S\)
produces electromagnetic wave propagation at speed \(c\), with the
physical content encoded in parameter ratios rather than absolute
values. The curl structure of Faraday's and Ampère's laws emerges from
projecting a depth-2 recursion (one hidden symbolic dimension) into
observable space. Five independent arguments --- MED node bounds, curl
algebra closure, Möbius embedding, orbital stability, and quaternion
uniqueness --- each independently require exactly three spatial
dimensions.

Charge is a topological winding number, quantised by the
single-valuedness of the phase field. Fractional quark charges
(\(\pm 1/3\), \(\pm 2/3\)) follow from the MED constraint that nodes
\(\leq 3\), creating three-fold internal substructure within each
topological defect.

As a speculative extension, clearly labelled: gravity may arise from the
same SEC wave equation through symmetric (divergence) projection, at
Fibonacci depth \(183 = F_7^2 + F_7 + 1\), giving
\(F_{183} \approx 10^{38}\) --- the observed
electromagnetic-to-gravitational hierarchy ratio. The LIGO/Virgo
measurement \(|c_\text{GW} - c_\text{EM}|/c < 3 \times 10^{-15}\) is
consistent with both forces sharing the same wave equation.

\textbf{Keywords}: Maxwell's equations, information geometry, curl
projection, spatial dimensions, MED bounds, charge quantisation,
Fibonacci structure, PAC conservation, Dawn Field Theory

\begin{center}\rule{0.5\linewidth}{0.5pt}\end{center}

\subsection{§1. What This Paper Does}\label{what-this-paper-does}

Classical electromagnetism is built on four equations discovered
empirically over the course of a century --- by Coulomb, Ampère,
Faraday, and Maxwell. These equations have a specific mathematical
structure: they use the curl operator, they apply in three spatial
dimensions, they involve an inverse-square force law, and they contain
two constants (\(\epsilon_0\), \(\mu_0\)) whose product gives \(1/c^2\).

None of this is explained by the equations themselves. Why curl and not
some other differential structure? Why three dimensions? Why
inverse-square? This paper presents arguments that these features follow
from three constraints established in Papers 1--4:

\begin{enumerate}
\def\labelenumi{\arabic{enumi}.}
\tightlist
\item
  \textbf{PAC conservation}: \(\Psi(k) = \Psi(k+1) + \Psi(k+2)\) --- the
  Fibonacci recursion
\item
  \textbf{MED bounds}: depth \(\leq 2\), nodes \(\leq 3\) --- the
  maximum complexity of emergent macroscopic patterns
\item
  \textbf{SEC dynamics}:
  \(\partial S/\partial t = \alpha\nabla I - \beta\nabla H\) ---
  information-entropy gradient flow
\end{enumerate}

The derivation follows a hierarchy:

\begin{verbatim}
Level 0 (PAC):     f(Parent) = Σf(Children), Ξ ≈ 1.0571
        ↓ SEC collapse (depth = 2)
Level 1 (Gauge):   sin²θ_W = F₄/F₇ = 3/13
        ↓ projection to 3+1D (MED: nodes ≤ 3)
Level 2 (Maxwell): E ↔ B curl structure, c = 1/√(ε₀μ₀)
\end{verbatim}

Each level is a projection of the one above. Maxwell's equations are
what the PAC recursion looks like when viewed through the MED filter at
depth 2 in 3 spatial dimensions.

\begin{center}\rule{0.5\linewidth}{0.5pt}\end{center}

\subsection{§2. The SEC Wave Equation}\label{the-sec-wave-equation}

\subsubsection{§2.1 From Gradient Flow to Wave
Propagation}\label{from-gradient-flow-to-wave-propagation}

The SEC (Symbolic Entropy Collapse) dynamic describes the evolution of
structure \(S\) under competing information and entropy gradients:

\[\frac{\partial S}{\partial t} = \alpha\nabla I - \beta\nabla H\]

where \(I\) is the information field, \(H\) is the entropy field, and
\(\alpha\), \(\beta\) are coupling coefficients. When \(I\) and \(H\)
are themselves functions of \(S\) through constitutive relations
(\(I = \gamma S\), \(H = \delta S\)), the first-order equation becomes
second-order:

\[\frac{\partial^2 S}{\partial t^2} = (\alpha\gamma + \beta\delta)\nabla^2 S\]

This is the wave equation with propagation speed:

\[c_\text{SEC}^2 = \alpha\gamma + \beta\delta\]

\subsubsection{§2.2 Parameter Models}\label{parameter-models}

The wave speed \(c\) is determined by the sum
\(\alpha\gamma + \beta\delta\), not by the individual parameters. Five
parameter models were tested, each encoding a different hypothesis about
the internal structure {[}1{]}:

\begin{longtable}[]{@{}
  >{\raggedright\arraybackslash}p{(\linewidth - 4\tabcolsep) * \real{0.2188}}
  >{\raggedright\arraybackslash}p{(\linewidth - 4\tabcolsep) * \real{0.3438}}
  >{\raggedright\arraybackslash}p{(\linewidth - 4\tabcolsep) * \real{0.4375}}@{}}
\toprule\noalign{}
\begin{minipage}[b]{\linewidth}\raggedright
Model
\end{minipage} & \begin{minipage}[b]{\linewidth}\raggedright
Hypothesis
\end{minipage} & \begin{minipage}[b]{\linewidth}\raggedright
Key Relation
\end{minipage} \\
\midrule\noalign{}
\endhead
\bottomrule\noalign{}
\endlastfoot
Symmetric & \(\alpha = \beta\), \(\gamma = \delta\) &
\(c^2 = 2\alpha\gamma\) \\
Ξ-balanced & \(\alpha/\beta = \Xi \approx 1.0571\) &
\(c^2 = \beta^2(\Xi + 1)\) \\
φ-structured & \(\alpha/\gamma = \varphi\), \(\beta/\delta = 1/\varphi\)
& \(c^2 = \gamma^2(\varphi + 1/\varphi)\) \\
Fibonacci-nested & \(\alpha/\beta = F_{10}/F_7 = 55/13\) &
\(c^2 = v_0^2(55/13 + 1)\) \\
Ξ-mean & \(\alpha/\beta = \sqrt{\Xi \cdot \Xi_\text{min}}\) &
\(c^2 = v_0^2(\Xi_\text{mean} + 1)\) \\
\end{longtable}

All five models produce \(c\) exactly --- by construction. The physical
content is not in the absolute values of
\(\alpha, \beta, \gamma, \delta\) but in their ratios, which encode the
balance between information accumulation and entropy dissipation. This
is consistent with the result from Paper 2 {[}2{]}: the balance constant
Ξ appears at the ratio level, not the absolute level.

The speed of light is not a free parameter in this framework. It is the
propagation speed of disturbances in the information-entropy field,
fixed by whatever internal structure sets the ratios.

\textbf{Script}: \texttt{exp\_01\_sec\_wave\_speed.py}

\begin{center}\rule{0.5\linewidth}{0.5pt}\end{center}

\subsection{§3. Why Three Spatial
Dimensions}\label{why-three-spatial-dimensions}

\subsubsection{§3.1 The Question}\label{the-question}

The Standard Model is formulated in 3+1 dimensions. String theory
requires 10. M-theory requires 11. Loop quantum gravity operates in 4.
None of these theories explains why the observed macroscopic universe
has exactly three large spatial dimensions. The number 3 is an input.

Here we present five independent arguments, each arriving at \(D = 3\)
from different starting points. No single argument is conclusive. Their
convergence is the result.

\subsubsection{§3.2 Path 1: MED Node Bound}\label{path-1-med-node-bound}

The MED (Macro Emergence Dynamics) constraint, discovered empirically in
Navier-Stokes simulations (see §6 and the theory experiments
repository), states that all complex flows converge to symbolic patterns
with:

\[\text{depth}(\sigma) \leq 2, \quad \text{nodes}(\sigma) \leq 3\]

across 1,000+ simulations spanning Reynolds numbers from 10 to 50,000.
If each spatial axis corresponds to an independent node in the emergent
pattern, the node bound directly gives:

\[D \leq 3\]

A fourth spatial dimension would require a fourth independent node,
violating the MED bound. This is not a mathematical proof that \(D > 3\)
is impossible --- it is an empirical observation that macroscopic
emergent structures do not sustain more than three independent
directions.

\subsubsection{§3.3 Path 2: Curl Algebra
Closure}\label{path-2-curl-algebra-closure}

The curl operator in \(n\) dimensions produces \(n(n-1)/2\) components
(the dimension of the exterior algebra \(\Lambda^2(\mathbb{R}^n)\)). For
the curl to map vectors to vectors (a prerequisite for Maxwell's
equations), we need:

\[\frac{n(n-1)}{2} = n\]

Solving: \(n^2 - 3n = 0\), giving \(n(n-3) = 0\), so \(n = 0\) or
\(n = 3\).

Only in three dimensions does the curl of a vector field produce another
vector field. In 2D, the curl produces a scalar. In 4D, it produces a
6-component object (an antisymmetric tensor, not a vector). Maxwell's
equations in their familiar form --- where \(\mathbf{E}\) and
\(\mathbf{B}\) are both 3-vectors related by curl --- exist only in 3D.

This is not new mathematics. It is the Hodge duality
\(\star: \Lambda^2(\mathbb{R}^3) \to \Lambda^1(\mathbb{R}^3)\), which is
an isomorphism only when \(\binom{n}{2} = n\). What is new is connecting
it to the MED bound: MED says nodes \(\leq 3\), and curl closure says
\(n = 3\). These are independent constraints that agree.

\subsubsection{§3.4 Path 3: Möbius
Embedding}\label{path-3-muxf6bius-embedding}

The PAC recursion operates on a Möbius-like topology (Paper 1 {[}1{]}):
the pre-field structure is non-orientable, with \(4\pi\) phase recovery.
A non-orientable 2-surface (Möbius strip, Klein bottle) cannot be
embedded without self-intersection in fewer than 3 dimensions.

This is a topological necessity: if the pre-field has Möbius topology,
the minimum observable spatial dimension is 3.

\textbf{Numerical verification}: A Möbius strip parameterised on a unit
square uses a non-trivial \(z\)-coordinate range of exactly 1.0 ---
confirming that the third dimension is structurally required, not
vestigial.

\subsubsection{§3.5 Path 4: Orbital Stability (Bertrand's
Theorem)}\label{path-4-orbital-stability-bertrands-theorem}

Gauss's law in \(n\) dimensions gives:

\[E \propto \frac{1}{r^{n-1}}\]

Bertrand's theorem {[}6{]} proves that closed, stable orbits under a
central force exist only for \(F \propto 1/r^2\) (Kepler) and
\(F \propto r\) (harmonic). The inverse-square law requires
\(n - 1 = 2\), i.e., \(n = 3\).

In \(D = 4\), orbits spiral inward; in \(D = 2\), the force is
logarithmic and orbits are marginally stable. Only \(D = 3\) supports
the stable planetary orbits, atomic structure, and bound states that
constitute observable matter.

\subsubsection{§3.6 Path 5: Quaternion
Uniqueness}\label{path-5-quaternion-uniqueness}

The rotation group \(SO(3)\) is locally isomorphic to
\(SU(2)/\mathbb{Z}_2\), the quaternion group modulo sign. There are
exactly four normed division algebras over \(\mathbb{R}\) (Hurwitz's
theorem {[}7{]}):

\begin{longtable}[]{@{}lll@{}}
\toprule\noalign{}
Algebra & Dimension & Rotation Group \\
\midrule\noalign{}
\endhead
\bottomrule\noalign{}
\endlastfoot
\(\mathbb{R}\) & 1 & trivial \\
\(\mathbb{C}\) & 2 & \(SO(2)\) \\
\(\mathbb{H}\) & 4 (3 imaginary) & \(SO(3)\) \\
\(\mathbb{O}\) & 8 (7 imaginary) & \(G_2\) \\
\end{longtable}

Three-dimensional rotations are uniquely associated with the quaternions
--- the only non-commutative, associative division algebra. The
octonions (\(\mathbb{O}\)) are non-associative and do not support a
conventional rotation group. If physical rotations must be associative
(which composition of spatial rotations requires), \(D = 3\) is the
unique answer.

\subsubsection{§3.7 Convergence}\label{convergence}

\begin{longtable}[]{@{}
  >{\raggedright\arraybackslash}p{(\linewidth - 6\tabcolsep) * \real{0.1463}}
  >{\raggedright\arraybackslash}p{(\linewidth - 6\tabcolsep) * \real{0.3659}}
  >{\raggedright\arraybackslash}p{(\linewidth - 6\tabcolsep) * \real{0.2927}}
  >{\raggedright\arraybackslash}p{(\linewidth - 6\tabcolsep) * \real{0.1951}}@{}}
\toprule\noalign{}
\begin{minipage}[b]{\linewidth}\raggedright
Path
\end{minipage} & \begin{minipage}[b]{\linewidth}\raggedright
Starting Point
\end{minipage} & \begin{minipage}[b]{\linewidth}\raggedright
Constraint
\end{minipage} & \begin{minipage}[b]{\linewidth}\raggedright
Result
\end{minipage} \\
\midrule\noalign{}
\endhead
\bottomrule\noalign{}
\endlastfoot
1 & MED bounds (empirical) & nodes \(\leq 3\) & \(D \leq 3\) \\
2 & Curl algebra (mathematical) & \(n(n-1)/2 = n\) & \(D = 3\) \\
3 & Möbius embedding (topological) & non-orientable surface &
\(D \geq 3\) \\
4 & Orbital stability (dynamical) & Bertrand's theorem & \(D = 3\) \\
5 & Division algebras (algebraic) & associative rotations & \(D = 3\) \\
\end{longtable}

Paths 1 and 3 provide upper and lower bounds: \(3 \leq D \leq 3\). Paths
2, 4, and 5 independently select \(D = 3\) from algebraic, dynamical,
and algebraic-topological requirements. No other integer satisfies all
five simultaneously.

\textbf{Script}: \texttt{exp\_05\_3d\_necessity.py}

\begin{center}\rule{0.5\linewidth}{0.5pt}\end{center}

\subsection{§4. Curl from Depth-2
Projection}\label{curl-from-depth-2-projection}

\subsubsection{§4.1 The Hidden Dimension}\label{the-hidden-dimension}

The SEC field \(S\) evolves in \(d + 1\) dimensions: \(d\) observable
spatial dimensions plus one symbolic recursion dimension. The MED
constraint fixes depth \(\leq 2\), which means this recursion dimension
is not directly observable --- it is the ``depth'' of the emergent
pattern, not a spatial direction.

When a \((d+1)\)-dimensional gradient is projected to \(d\) observable
dimensions, the component along the hidden dimension becomes a source of
curl:

\[\nabla_{d+1}\Phi \xrightarrow{\text{project out } z_\text{hidden}} (\nabla_d \Phi, \partial\Phi/\partial z_\text{hidden})\]

The hidden component \(\partial\Phi/\partial z_\text{hidden}\) manifests
as rotation in the observable \(d\)-dimensional space. This is the
origin of the magnetic field: it is the projection artifact of a
gradient in the recursion dimension.

\subsubsection{§4.2 Faraday's Law}\label{faradays-law}

Faraday's law of induction:

\[\nabla \times \mathbf{E} = -\frac{\partial \mathbf{B}}{\partial t}\]

In the SEC framework: when the magnetic field \(\mathbf{B}\) (=
projected recursion gradient) changes in time, the information vorticity
must compensate to conserve the total SEC flux. The electric field
\(\mathbf{E}\) acquires curl to balance the changing \(\mathbf{B}\).

Numerical verification: for sinusoidal \(\mathbf{E}\) and \(\mathbf{B}\)
fields satisfying the SEC evolution,
\(|\nabla \times \mathbf{E} + \partial\mathbf{B}/\partial t| \approx 10^{-16}\)
(machine precision).

\subsubsection{§4.3 Depth 1 vs Depth 2}\label{depth-1-vs-depth-2}

\begin{longtable}[]{@{}
  >{\raggedright\arraybackslash}p{(\linewidth - 4\tabcolsep) * \real{0.2593}}
  >{\raggedright\arraybackslash}p{(\linewidth - 4\tabcolsep) * \real{0.4074}}
  >{\raggedright\arraybackslash}p{(\linewidth - 4\tabcolsep) * \real{0.3333}}@{}}
\toprule\noalign{}
\begin{minipage}[b]{\linewidth}\raggedright
Depth
\end{minipage} & \begin{minipage}[b]{\linewidth}\raggedright
Structure
\end{minipage} & \begin{minipage}[b]{\linewidth}\raggedright
Physics
\end{minipage} \\
\midrule\noalign{}
\endhead
\bottomrule\noalign{}
\endlastfoot
1 & Only gradients; no cross-derivatives & Electrostatics only (Coulomb,
no magnetism) \\
2 & Cross-derivatives exist; curl emerges & Full electromagnetism (E, B
coupled) \\
\(\geq 3\) & MED says systems converge back to \(\leq 2\) & Unstable;
decays to depth 2 \\
\end{longtable}

The depth-2 MED bound is not just a limit --- it is an attractor.
Systems that transiently reach depth 3 collapse back to depth 2. This is
why electromagnetism is the only long-range force with curl structure:
it saturates the MED depth bound.

\subsubsection{§4.4 MED Saturation and
Dimensionality}\label{med-saturation-and-dimensionality}

The total effective depth of a system with \(d\) spatial dimensions:

\[d_\text{total} = d_\text{physical} + d_\text{symbolic} = d + 1\]

For MED saturation (\$d\_\text{total} = \$ depth bound \(\times 2\)):

\begin{longtable}[]{@{}llll@{}}
\toprule\noalign{}
\(d\) & \(d_\text{total}\) & Effective depth & MED saturated? \\
\midrule\noalign{}
\endhead
\bottomrule\noalign{}
\endlastfoot
2 & 3 & 1.5 & No \\
3 & 4 & 2.0 & \textbf{Yes} \\
4 & 5 & 2.5 & Exceeds bound \\
\end{longtable}

Three dimensions is the unique value that exactly saturates the MED
depth-2 bound. The She-Lévêque intermittency formula with full cascade
structure (\(k = 9\); Paper 4 {[}9{]}, §9) requires this saturation ---
in 2D, the cascade is modified (\(k = 4\)) because MED is not saturated.

\textbf{Script}: \texttt{exp\_03\_curl\_projection.py}

\begin{center}\rule{0.5\linewidth}{0.5pt}\end{center}

\subsection{§5. Charge as Winding
Number}\label{charge-as-winding-number}

\subsubsection{§5.1 Topological
Quantisation}\label{topological-quantisation}

In the SEC field, a point defect creates a phase singularity. The phase
\(\theta\) of the SEC order parameter must be single-valued after
traversing any closed loop around the defect:

\[n = \frac{1}{2\pi}\oint d\theta = \text{integer}\]

This is the winding number. It is topologically protected: continuous
deformations of the field cannot change \(n\). The winding number is the
electric charge, measured in units of \(e\).

\subsubsection{§5.2 Coulomb's Law}\label{coulombs-law}

A point defect with winding number \(n\) in 3D creates a phase gradient
that falls off as \(1/r^2\) (Gauss's law applied to the solid angle
subtended by the defect):

\[\mathbf{E} = \frac{n}{4\pi\epsilon_0 r^2}\hat{r}\]

Numerical verification: the field of a unit winding defect falls off as
\(r^{-2.0000}\) to four decimal places.

\subsubsection{§5.3 Pair Creation and PAC
Conservation}\label{pair-creation-and-pac-conservation}

Pair creation (\(\gamma \to e^+ e^-\)) requires total winding number
zero:

\[n_\text{total} = n_{e^-} + n_{e^+} = (+1) + (-1) = 0\]

This is PAC conservation applied to topological charge: the parent
(photon, \(n = 0\)) equals the sum of children (\(n = +1\) and
\(n = -1\)).

\subsubsection{§5.4 Fractional Charges and the MED
Constraint}\label{fractional-charges-and-the-med-constraint}

Quarks carry fractional charges (\(+2/3\), \(-1/3\)) {[}8{]}. Within the
topological framework, this requires sub-defect structure: a single
winding defect contains internal nodes.

The MED constraint (nodes \(\leq 3\)) limits this internal structure to
at most three sub-nodes. Each sub-node independently carries winding
\(+1/3\), \(-1/3\), or \(0\). The particle's charge is the sum of its
sub-node windings:

\begin{itemize}
\tightlist
\item
  \textbf{Up-type quarks} (\(q = +2/3\)): two sub-nodes carry \(+1/3\)
  each, one carries \(0\)
\item
  \textbf{Down-type quarks} (\(q = -1/3\)): one sub-node carries
  \(-1/3\), two carry \(0\)
\item
  \textbf{Charged leptons} (\(q = -1\)): all three sub-nodes carry
  \(-1/3\)
\end{itemize}

Colour confinement follows: observable states must have integer total
winding. A baryon combines three quarks (\(uud\):
\(+2/3 + 2/3 - 1/3 = +1\)), and a meson combines quark and antiquark
(\(u\bar{d}\): \(+2/3 + 1/3 = +1\)). Free quarks, with fractional total
winding, cannot be isolated --- the MED tree is not closed.

The observed quark charges are the only assignments consistent with:

\begin{enumerate}
\def\labelenumi{\arabic{enumi}.}
\tightlist
\item
  Total winding = integer for colour-singlet states
\item
  Internal nodes \(\leq 3\) (MED bound)
\item
  PAC conservation at each node
\end{enumerate}

This does not derive the full quark charge assignments from first
principles --- it constrains them to a small set. The specific
assignment (\(+2/3\) for up-type, \(-1/3\) for down-type) may follow
from the Fibonacci gauge structure (Paper 4 {[}9{]}), but this
connection is not yet established.

\textbf{Script}: \texttt{exp\_02\_charge\_quantization.py}

\begin{center}\rule{0.5\linewidth}{0.5pt}\end{center}

\subsection{§6. The SEC--Navier-Stokes
Equivalence}\label{the-secnavier-stokes-equivalence}

\subsubsection{§6.1 Structural Mapping}\label{structural-mapping}

The SEC dynamics:

\[\frac{\partial F}{\partial t} = -\alpha\nabla H(F) + \beta\mathcal{D}(F) + \gamma\mathcal{M}(F, t)\]

maps term-by-term onto the Navier-Stokes equations:

\begin{longtable}[]{@{}
  >{\raggedright\arraybackslash}p{(\linewidth - 4\tabcolsep) * \real{0.3571}}
  >{\raggedright\arraybackslash}p{(\linewidth - 4\tabcolsep) * \real{0.3214}}
  >{\raggedright\arraybackslash}p{(\linewidth - 4\tabcolsep) * \real{0.3214}}@{}}
\toprule\noalign{}
\begin{minipage}[b]{\linewidth}\raggedright
SEC Term
\end{minipage} & \begin{minipage}[b]{\linewidth}\raggedright
NS Term
\end{minipage} & \begin{minipage}[b]{\linewidth}\raggedright
Physics
\end{minipage} \\
\midrule\noalign{}
\endhead
\bottomrule\noalign{}
\endlastfoot
\(-\alpha\nabla H\) & \(-\nabla p\) & Pressure gradient \\
\(\beta\mathcal{D}\) & \(\nu\nabla^2 \mathbf{u}\) & Viscous diffusion \\
\(\gamma\mathcal{M}\) & \((\mathbf{u} \cdot \nabla)\mathbf{u}\) &
Advective nonlinearity \\
\end{longtable}

This is a structural equivalence, not an analogy. The SEC equation with
appropriate constitutive relations \emph{is} the Navier-Stokes equation.
The balance operator \(\Xi \approx 1.0571\) prevents finite-time blowup
--- the same role it plays in regulating coupling constants (Paper 2
{[}2{]}).

\subsubsection{§6.2 Regularity from MED
Bounds}\label{regularity-from-med-bounds}

The Navier-Stokes regularity problem (Clay Millennium Prize) asks
whether smooth initial data always lead to smooth solutions. The MED
bounds provide a route to regularity:

\begin{itemize}
\tightlist
\item
  \textbf{Depth \(\leq 2\)} implies gradient control:
  \(\|\nabla\mathbf{u}(t)\|_\infty \leq C_\text{global}\) for all \(t\)
\item
  \textbf{Nodes \(\leq 3\)} implies energy control:
  \(\|\mathbf{u}(t)\|^2_{L^2} \leq E_\text{total}\) for all \(t\)
\end{itemize}

Together, these prevent the formation of finite-time singularities. Over
10,000 pattern transitions were tested with 100\% Landauer compliance
and energy conservation error \(< 10^{-12}\).

This is not a proof of Navier-Stokes regularity --- it is an empirical
observation that MED-bounded SEC dynamics do not blow up, plus a
structural argument for why. The connection to a formal proof would
require showing that the MED bounds are not just observed but necessary
consequences of the SEC structure.

\subsubsection{§6.3 Turbulence Constants as Fibonacci
Structure}\label{turbulence-constants-as-fibonacci-structure}

Paper 4 {[}9{]} (§9) established that the She-Lévêque intermittency
parameters are Fibonacci:

\[\zeta_p = \frac{p}{(F_4)^2} + F_3\left[1 - \left(\frac{F_3}{F_4}\right)^{p/F_4}\right] = \frac{p}{9} + 2\left[1 - \left(\frac{2}{3}\right)^{p/3}\right]\]

The dimensional formula \(k(d) = d \times F_{d+1}\) connects turbulence
to spatial dimensionality through the Fibonacci sequence:

\begin{longtable}[]{@{}llll@{}}
\toprule\noalign{}
\(d\) & \(k\) & Fibonacci & Verification \\
\midrule\noalign{}
\endhead
\bottomrule\noalign{}
\endlastfoot
2 & 4 & \(2 \times F_3 = 2 \times 2\) & 2D enstrophy cascade {[}3{]} \\
3 & 9 & \(3 \times F_4 = 3 \times 3\) & She-Lévêque (3D) {[}3{]} \\
4 & 20 & \(4 \times F_5 = 4 \times 5\) & \textbf{Prediction} \\
\end{longtable}

The cascade fraction \(\beta = F_3/F_4 = 2/3\) is the same ratio that
gives the Koide formula (Paper 4 {[}9{]}, §5). The Kolmogorov spectral
exponent \(5/3 = F_5/F_4\) is a ratio of consecutive Fibonacci numbers
--- a factual observation, not a new derivation of the Kolmogorov
spectrum, which follows from dimensional analysis (Kolmogorov, 1941
{[}4{]}).

The connection to electromagnetism: MED saturation at \(d = 3\) is the
same condition that enables the full She-Lévêque cascade (\(k = 9\)).
The turbulence and electromagnetic structures are both consequences of
depth-2 recursion in 3 spatial dimensions.

\begin{center}\rule{0.5\linewidth}{0.5pt}\end{center}

\subsection{§7. Mersenne Dimensions and Vacuum
Regularisation}\label{mersenne-dimensions-and-vacuum-regularisation}

\subsubsection{§7.1 Casimir 240}\label{casimir-240}

The Casimir force between parallel conducting plates:

\[\frac{F}{A} = -\frac{\pi^2 \hbar c}{240\,a^4}\]

The factor 240 arises from zeta function regularisation of the vacuum
mode sum. In Fibonacci arithmetic:

\[240 = F_3 \times F_4 \times F_5 \times F_6 = 2 \times 3 \times 5 \times 8\]

Four consecutive Fibonacci numbers. The sub-factor
\(120 = F_4 \times F_5 \times F_6\) (three consecutive) corresponds to
\(|\zeta(-3)| = 1/120\); the additional \(F_3 = 2\) enters through the
two-plate geometry.

\subsubsection{§7.2 Mersenne Dimensions}\label{mersenne-dimensions}

Fibonacci product structure in zeta function regularisation denominators
appears only at Mersenne dimensions \(d = 2^n - 1\):

\begin{longtable}[]{@{}llll@{}}
\toprule\noalign{}
\(d\) & \(\zeta(-d)\) denominator & Fibonacci product? & Mersenne? \\
\midrule\noalign{}
\endhead
\bottomrule\noalign{}
\endlastfoot
1 & 12 & \(F_3^2 \times F_4\) & \(2^1 - 1\) ✓ \\
3 & 120 & \(F_4 \times F_5 \times F_6\) & \(2^2 - 1\) ✓ \\
5 & 252 & No (factor 7, not Fibonacci) & ✗ \\
7 & 240 & \(F_3 \times F_4 \times F_5 \times F_6\) & \(2^3 - 1\) ✓ \\
9 & 132 & No (factor 11, not Fibonacci) & ✗ \\
\end{longtable}

Non-Mersenne dimensions have non-Fibonacci prime factors in their
regularisation denominators. The Mersenne dimensions that do have
Fibonacci structure are precisely those relevant to physical theories:

\begin{itemize}
\tightlist
\item
  \(d = 1\) (\(2^1 - 1\)): string worldsheet
\item
  \(d = 3\) (\(2^2 - 1\)): physical space
\item
  \(d = 7\) (\(2^3 - 1\)): M-theory extra dimensions
\end{itemize}

\subsubsection{§7.3 Falsifiable Extension}\label{falsifiable-extension}

If this pattern holds, the next Mersenne dimension \(d = 15\)
(\(2^4 - 1\)) should also have Fibonacci-structured regularisation. The
denominator of \(\zeta(-15)\) is
\(8160 = 2^5 \times 3 \times 5 \times 17\). The prime factor 17 is not a
Fibonacci number.

This falsifies the simple Mersenne-Fibonacci extension at \(d = 15\).
The pattern that holds for \(d = 1, 3, 7\) does not continue to
\(d = 15\). Either the correspondence is limited to the first three
Mersenne primes (perhaps connected to the existence of normed division
algebras in dimensions 1, 2, 4, 8 --- i.e., \(\mathbb{R}\),
\(\mathbb{C}\), \(\mathbb{H}\), \(\mathbb{O}\)), or the relevant
structure is a sub-factor of the denominator rather than the full
factorisation. We record this as an honest negative result --- the
pattern is interesting for \(d \leq 7\) but does not generalise
trivially.

\textbf{Scripts}: \texttt{exp\_15\_casimir\_sec.py},
\texttt{exp\_16\_mersenne\_verification.py}

\begin{center}\rule{0.5\linewidth}{0.5pt}\end{center}

\subsection{§8. Speculative Extension:
Gravity}\label{speculative-extension-gravity}

\emph{The content of this section is speculative. It is clearly labelled
as such and is not part of the core derivation chain (§§2--7). It is
included because the framework suggests it and because it makes a
testable order-of-magnitude prediction.}

\subsubsection{§8.1 Projection Duality}\label{projection-duality}

Any rank-2 tensor decomposes into symmetric and antisymmetric parts:

\[T = S + A, \quad S = \frac{T + T^\top}{2}, \quad A = \frac{T - T^\top}{2}\]

\begin{longtable}[]{@{}
  >{\raggedright\arraybackslash}p{(\linewidth - 6\tabcolsep) * \real{0.1739}}
  >{\raggedright\arraybackslash}p{(\linewidth - 6\tabcolsep) * \real{0.2754}}
  >{\raggedright\arraybackslash}p{(\linewidth - 6\tabcolsep) * \real{0.3188}}
  >{\raggedright\arraybackslash}p{(\linewidth - 6\tabcolsep) * \real{0.2319}}@{}}
\toprule\noalign{}
\begin{minipage}[b]{\linewidth}\raggedright
Projection
\end{minipage} & \begin{minipage}[b]{\linewidth}\raggedright
Degrees of Freedom
\end{minipage} & \begin{minipage}[b]{\linewidth}\raggedright
Differential Operator
\end{minipage} & \begin{minipage}[b]{\linewidth}\raggedright
Physical Theory
\end{minipage} \\
\midrule\noalign{}
\endhead
\bottomrule\noalign{}
\endlastfoot
Antisymmetric \(A\) & 3 (in 3D) & Curl (\(\nabla \times\)) &
Electromagnetism (\(F_{\mu\nu}\)) \\
Symmetric \(S\) & 6 (in 3D) & Divergence (\(\nabla \cdot\)) & Gravity
(\(h_{\mu\nu}\)) \\
\end{longtable}

The hypothesis: electromagnetism and gravity are the two natural
projections of the same SEC tensor field. EM uses the antisymmetric
projection (phase → curl → charge = winding number, integer-quantised).
Gravity uses the symmetric projection (amplitude → divergence → mass =
resonance frequency, continuously valued).

\subsubsection{§8.2 The Hierarchy Depth}\label{the-hierarchy-depth}

If EM operates at Fibonacci depth \(F_7 = 13\) (Paper 4 {[}9{]}, §3.2),
gravity operates at depth:

\[183 = F_7^2 + F_7 + 1 = 169 + 13 + 1\]

The polynomial \(N^2 + N + 1\) counts elements of the projective plane
\(PG(2, N)\). The Fibonacci number at this depth:

\[F_{183} \approx 1.27 \times 10^{38}\]

The measured electromagnetic-to-gravitational hierarchy ratio:

\[\frac{F_\text{EM}}{F_\text{grav}} = \frac{e^2/(4\pi\epsilon_0)}{Gm_p^2} \approx 1.24 \times 10^{38}\]

Same order of magnitude. As discussed in Paper 4 {[}9{]} (§11), this is
suggestive but not a precision prediction --- multiple nearby Fibonacci
indices give values within an order of magnitude of \(10^{38}\).

\subsubsection{§8.3 Same Wave Equation}\label{same-wave-equation}

If both EM and gravity originate from the SEC wave equation, they should
propagate at the same speed. The LIGO/Virgo observation of GW170817
{[}5{]} confirms:

\[\frac{|c_\text{GW} - c_\text{EM}|}{c} < 3 \times 10^{-15}\]

This is consistent with both forces sharing
\(c_\text{SEC}^2 = \alpha\gamma + \beta\delta\), with different
projections at different hierarchy depths.

\subsubsection{§8.4 Cosmic Structure from Local PAC
Gravity}\label{cosmic-structure-from-local-pac-gravity}

N-body simulations using PAC-local gravity (\(F \propto e^{-r/r_0}/r\),
short-range) produce cosmic web structure: - Void fraction: 50--89\% -
Power spectrum slope: \(n = -1.73\) (\(R^2 = 0.57\)) - 85\% match to
observed cosmic matter power spectrum statistics

Standard cosmology requires Newtonian \(1/r^2\) at all scales. The PAC
gravity simulation suggests that local exponential interactions,
iterated through PAC conservation, may reproduce large-scale structure
without long-range gravity. This is preliminary and does not replace the
standard model of cosmology --- it is a numerical observation that the
statistical signatures are similar.

\textbf{Scripts}: \texttt{exp\_09\_pac\_web.py},
\texttt{exp\_12\_power\_spectrum.py} (gravity\_from\_maxwell\_pac)

\begin{center}\rule{0.5\linewidth}{0.5pt}\end{center}

\subsection{§9. The Balance Constant in Wave
Dynamics}\label{the-balance-constant-in-wave-dynamics}

\subsubsection{§9.1 Ξ from SEC Dynamics}\label{ux3be-from-sec-dynamics}

The balance constant \(\Xi = 1 + \pi/55 \approx 1.0571\) (Paper 2
{[}2{]}) was discovered empirically in the Navier-Stokes symbolic
engine. Its derivation from SEC dynamics:

The SEC collapse rate per recursion level has two components:

\begin{itemize}
\tightlist
\item
  \textbf{Within-level}:
  \(\Delta_\text{within} = 2\sqrt{r(1-r)} - 1 = -0.0283\) per level (at
  \(r = 1/\varphi\))
\item
  \textbf{Cross-level}: \(\Delta_\text{cross} = +0.0854\) per level
  (network interference)
\item
  \textbf{Net}:
  \(\Delta_\text{net} = -0.0283 + 0.0854 = 0.0571 = \pi/55\) per level
\end{itemize}

At depth \(F_{10} = 55\):

\[55 \times \frac{\pi}{55} = \pi\]

One Möbius half-twist. The balance constant \(\Xi - 1 = \pi/55\) is the
per-level accumulation rate that produces exactly \(\pi\) radians of
phase shift over the full hierarchy depth. This closes the explanatory
loop: \(F_{10} = 55\) is not arbitrary --- it is the depth at which SEC
dynamics complete one half-twist of the Möbius pre-field topology.

\subsubsection{§9.2 Ξ as Blowup
Regulator}\label{ux3be-as-blowup-regulator}

In the SEC--Navier-Stokes equivalence (§6.1), the balance operator
\(\Xi\) prevents finite-time singularity formation. Its role is
analogous to the viscosity coefficient \(\nu\) but operates at the
symbolic level: it bounds the rate at which complexity can increase per
recursion step.

A/B testing across diverse dynamical systems shows that Ξ-correct
thresholds produce 1.48× faster convergence compared to arbitrary
thresholds, and Ξ-wrong thresholds produce 50.96× slower convergence.
Combined significance: \(p < 0.00001\).

\begin{center}\rule{0.5\linewidth}{0.5pt}\end{center}

\subsection{§10. The Complete Derivation
Chain}\label{the-complete-derivation-chain}

Assembling §§2--9:

\begin{longtable}[]{@{}
  >{\raggedright\arraybackslash}p{(\linewidth - 6\tabcolsep) * \real{0.2000}}
  >{\raggedright\arraybackslash}p{(\linewidth - 6\tabcolsep) * \real{0.2333}}
  >{\raggedright\arraybackslash}p{(\linewidth - 6\tabcolsep) * \real{0.2667}}
  >{\raggedright\arraybackslash}p{(\linewidth - 6\tabcolsep) * \real{0.3000}}@{}}
\toprule\noalign{}
\begin{minipage}[b]{\linewidth}\raggedright
Step
\end{minipage} & \begin{minipage}[b]{\linewidth}\raggedright
Input
\end{minipage} & \begin{minipage}[b]{\linewidth}\raggedright
Output
\end{minipage} & \begin{minipage}[b]{\linewidth}\raggedright
Section
\end{minipage} \\
\midrule\noalign{}
\endhead
\bottomrule\noalign{}
\endlastfoot
1 & PAC recursion \(\Psi(k) = \Psi(k+1) + \Psi(k+2)\) & Fibonacci
numbers, \(\varphi\) & Paper 1 {[}1{]} \\
2 & SEC dynamics on PAC fields & Wave equation:
\(c^2 = \alpha\gamma + \beta\delta\) & §2 \\
3 & MED: nodes \(\leq 3\) & \(D = 3\) spatial dimensions & §3 \\
4 & MED: depth \(\leq 2\) & Curl from hidden dimension projection &
§4 \\
5 & Phase defect topology & Charge = winding number (integer) & §5 \\
6 & MED: nodes \(\leq 3\) in sub-defects & Quark charges \(\pm 1/3\),
\(\pm 2/3\) & §5.4 \\
7 & SEC--NS equivalence & Turbulence parameters are Fibonacci & §6 \\
8 & Fibonacci regularisation & Casimir 240, Mersenne dimensions & §7 \\
9 & Depth 55 = \(F_{10}\) & \(\Xi - 1 = \pi/55\) (one half-twist) &
§9 \\
10 & \emph{Speculative}: symmetric projection & Gravity at depth 183 &
§8 \\
\end{longtable}

Steps 1--9 form the core derivation. Step 10 is clearly separated as
speculative.

The key connective tissue is the MED bound. Without it: - We have no
reason for \(D = 3\) (§3) - We have no curl structure (§4) - We have no
quark charge quantisation (§5.4) - We have no She-Lévêque saturation
(§6.3) - We have no NS regularity argument (§6.2)

MED is the single constraint that converts the abstract PAC/SEC
framework into the specific structure of classical physics.

\begin{center}\rule{0.5\linewidth}{0.5pt}\end{center}

\subsection{§11. What This Paper Does Not
Do}\label{what-this-paper-does-not-do}

This paper does not derive Maxwell's equations from axioms. A derivation
would start from the PAC recursion and, without reference to any
measured quantity, produce the four Maxwell equations with their
specific coupling constant. We do not achieve this.

What we show is that the \emph{structure} of Maxwell's equations ---
curl operations, three dimensions, inverse-square force, quantised
charge --- follows from information-theoretic constraints (PAC, MED,
SEC) that are independently motivated. The coupling constant \(\alpha\)
is expressed in Fibonacci arithmetic (Paper 4 {[}9{]}) but not derived
from first principles.

The speculative extension to gravity (§8) is even weaker: it provides an
order-of-magnitude prediction and a structural analogy (antisymmetric vs
symmetric projection), but no dynamical equation for gravity comparable
to Maxwell's equations for electromagnetism.

These limitations are stated not as caveats to be dismissed but as the
current status of the work. If the MED bounds can be derived from the
PAC recursion (rather than observed empirically), and if the hierarchy
depth formula \(n = F_7^2 + F_7 + 1\) can be justified structurally
(rather than noted post hoc), the framework would gain substantially.

\begin{center}\rule{0.5\linewidth}{0.5pt}\end{center}

\subsection{§12. Falsification
Conditions}\label{falsification-conditions}

\begin{enumerate}
\def\labelenumi{\arabic{enumi}.}
\item
  \textbf{MED bounds violated.} If a macroscopic emergent pattern with
  depth \(> 2\) or nodes \(> 3\) is found in any physical system and is
  demonstrably stable, the MED constraint fails and the \(D = 3\)
  argument (§3.2) collapses.
\item
  \textbf{Curl structure in non-3D.} If a self-consistent classical
  electrodynamics with vectors (not tensors) is constructed in
  \(D \neq 3\), the curl closure argument (§3.3) is invalidated. This is
  unlikely given the mathematics, but the physical claim could still be
  wrong if nature does not require vector-to-vector curl.
\item
  \textbf{\(k(4) \neq 20\).} If 4D turbulence simulations produce
  \(k \neq 20\), the formula \(k = d \times F_{d+1}\) fails and the
  Fibonacci turbulence structure loses predictive power.
\item
  \textbf{Mersenne pattern breaks at \(d = 15\).} The denominator of
  \(\zeta(-15)\) contains the non-Fibonacci prime 17 (§7.3). The
  Mersenne-Fibonacci correspondence holds for \(d = 1, 3, 7\) but does
  not extend to \(d = 15\), limiting the pattern to normed division
  algebra dimensions.
\item
  \textbf{\(c_\text{GW} \neq c_\text{EM}\) at higher precision.} If
  future observations resolve \(c_\text{GW} \neq c_\text{EM}\), the
  shared SEC wave equation (§8.3) is falsified.
\end{enumerate}

\begin{center}\rule{0.5\linewidth}{0.5pt}\end{center}

\subsection{§13. Connections to the
PACSeries}\label{connections-to-the-pacseries}

\begin{longtable}[]{@{}
  >{\raggedright\arraybackslash}p{(\linewidth - 2\tabcolsep) * \real{0.3889}}
  >{\raggedright\arraybackslash}p{(\linewidth - 2\tabcolsep) * \real{0.6111}}@{}}
\toprule\noalign{}
\begin{minipage}[b]{\linewidth}\raggedright
Paper
\end{minipage} & \begin{minipage}[b]{\linewidth}\raggedright
Connection
\end{minipage} \\
\midrule\noalign{}
\endhead
\bottomrule\noalign{}
\endlastfoot
Paper 1 {[}1{]}: Structure Cost of Erasure & SEC dynamics established;
Landauer structure cost gives \(\alpha\) its physical interpretation \\
Paper 2 {[}2{]}: Balance Constant & \(\Xi = 1 + \pi/55\) derived;
appears in SEC wave equation parameter ratios and as NS blowup
regulator \\
Paper 3 {[}10{]}: Feigenbaum Constants & \(F_{10} = 55\) in closed-form
Feigenbaum expressions; MED discovered in NS simulations \\
Paper 4 {[}9{]}: Standard Model Parameters & \(\alpha\) formula,
She-Lévêque \(k = 9\), Casimir 240, Mersenne dimensions, gauge closure
\(F_7 = 13\) \\
\textbf{Paper 5 (this paper)} & \textbf{Classical physics structure from
information geometry} \\
Paper 6: Computational Validation & PAC conservation observed in ML
systems --- the operational principle tested in artificial systems \\
\end{longtable}

The progression: Paper 1 establishes the cost. Papers 2--3 find the
constants. Paper 4 matches the Standard Model. Paper 5 shows why
classical physics has the structure it does. Paper 6 tests whether the
principle operates in systems we can build.

\begin{center}\rule{0.5\linewidth}{0.5pt}\end{center}

\subsection{§14. Summary}\label{summary}

\begin{longtable}[]{@{}
  >{\raggedright\arraybackslash}p{(\linewidth - 4\tabcolsep) * \real{0.3333}}
  >{\raggedright\arraybackslash}p{(\linewidth - 4\tabcolsep) * \real{0.3333}}
  >{\raggedright\arraybackslash}p{(\linewidth - 4\tabcolsep) * \real{0.3333}}@{}}
\toprule\noalign{}
\begin{minipage}[b]{\linewidth}\raggedright
Result
\end{minipage} & \begin{minipage}[b]{\linewidth}\raggedright
Source
\end{minipage} & \begin{minipage}[b]{\linewidth}\raggedright
Status
\end{minipage} \\
\midrule\noalign{}
\endhead
\bottomrule\noalign{}
\endlastfoot
SEC wave equation → \(c\) & §2 & Structural (exact by construction) \\
\(D = 3\) from 5 independent paths & §3 & Convergent (mathematical +
empirical) \\
Curl from depth-2 projection & §4 & Structural \\
Faraday's law to \(10^{-16}\) & §4.2 & Numerical \\
Charge = winding number & §5 & Topological \\
Coulomb \(r^{-2.0000}\) & §5.2 & Numerical \\
Quark charges from MED \(\leq 3\) & §5.4 & Constrained (not fully
derived) \\
SEC--NS equivalence & §6 & Structural mapping \\
MED → NS regularity & §6.2 & Empirical (\(10^4\) transitions,
\(10^{-12}\) conservation) \\
\(k = d \times F_{d+1}\) & §6.3 & Formula (verified \(d = 2, 3\);
prediction \(d = 4\)) \\
Casimir 240 = \(F_3 F_4 F_5 F_6\) & §7.1 & Exact identity \\
Mersenne-Fibonacci correspondence & §7.2 & Observed (\(d = 1, 3, 7\)) \\
\(\Xi - 1 = \pi/55\) derivation & §9 & Derived from SEC rates \\
Gravity at depth 183 & §8 & \textbf{Speculative} (order of magnitude) \\
\(c_\text{GW} = c_\text{EM}\) & §8.3 & Consistent with observation \\
\end{longtable}

\begin{center}\rule{0.5\linewidth}{0.5pt}\end{center}

\subsection{Acknowledgments}\label{acknowledgments}

\emph{(To be added.)}

\begin{center}\rule{0.5\linewidth}{0.5pt}\end{center}

\subsection{References}\label{references}

\begin{enumerate}
\def\labelenumi{\arabic{enumi}.}
\tightlist
\item
  Groom, P. (2026). ``The Structure Cost of Erasure.'' PACSeries Paper
  1. Dawn Field Institute.
\item
  Groom, P. (2026). ``The Balance Constant and Its Decomposition.''
  PACSeries Paper 2. Dawn Field Institute.
\item
  She, Z.-S. and Lévêque, E. (1994). ``Universal Scaling Laws in Fully
  Developed Turbulence.'' \emph{Phys. Rev.~Lett.}, 72(3), 336--339.
\item
  Kolmogorov, A. N. (1941). ``The local structure of turbulence in
  incompressible viscous fluid for very large Reynolds numbers.''
  \emph{Dokl. Akad. Nauk SSSR}, 30, 299--303.
\item
  Abbott, B. P. et al.~(2017). ``Gravitational Waves and Gamma-Rays from
  a Binary Neutron Star Merger: GW170817 and GRB 170817A.''
  \emph{Astrophys. J. Lett.}, 848, L13.
\item
  Bertrand, J. (1873). ``Théorème relatif au mouvement d'un point attiré
  vers un centre fixe.'' \emph{C. R. Acad. Sci. Paris}, 77, 849--853.
\item
  Hurwitz, A. (1898). ``Über die Composition der quadratischen Formen
  von beliebig vielen Variablen.'' \emph{Nachr. Ges. Wiss. Göttingen},
  309--316.
\item
  Particle Data Group (2022). ``Review of Particle Physics.''
  \emph{Prog. Theor. Exp. Phys.}, 2022(8), 083C01.
\item
  Groom, P. (2026). ``Standard Model Parameters from Fibonacci
  Arithmetic.'' PACSeries Paper 4. Dawn Field Institute.
\item
  Groom, P. (2026). ``Feigenbaum Constants from Fibonacci Arithmetic.''
  PACSeries Paper 3. Dawn Field Institute.
\end{enumerate}

\begin{center}\rule{0.5\linewidth}{0.5pt}\end{center}

\emph{All code, data, and experiment scripts for this paper and the full
PACSeries are publicly available at
\url{https://github.com/dawnfield-institute/dawn-field-theory}. See the
accompanying publication package README.md for reproduction
instructions.}

\end{document}
